% Programming with InitechBase, Chapter 12: Handling Errors (AN-16).
% Panels replayed (examples/p12-*.dot); sources samir_main.c (the message
% catalogue, "n  text"), cmd/proc.c (ON ERROR, proc_run), core/fn_builtins.c
% ERROR().  The program BADTOTAL.PRG of the sample diskette.
\chapter{Handling Errors}
\Gates{test-pubs-examples, test-interp-proc, test-samir-repl}
A command that cannot be carried out stops the program.  This chapter
describes the messages InitechBase displays, how to find the command at fault,
and how a program can record an error with \cmd{ON ERROR}.

\begin{ChapterCovers}
\item Error messages and their numbers
\item What happens to a program
\item \cmd{ON ERROR} and \fn{ERROR()}
\end{ChapterCovers}

\ATEHead{Preparing for This Chapter}
Start InitechBase with a fresh copy of the sample diskette.

\ATEHead{Error Messages}
When a command cannot be carried out, InitechBase displays a message: its
number, two blanks and its text.  Three common errors at the dot prompt:
\DotExample{examples/p12-errors}
\xmsg{9}{Data type mismatch.} means that an operator was given values of types
it cannot combine; \xmsg{12}{Variable not found.} that a name is neither a field
of the database in USE nor a memory variable; and the third that a function was
given an argument out of its range.  Chapter~7 of \doctitle{Using InitechBase}
lists every message with its cause and what to do.

\ATEHead{What Happens to a Program}
When a command of a program fails, the program stops at that command.
InitechBase displays the message and the dot prompt.  The memory variables the
program created remain, and the record pointer stays where it was; both help
you find the command at fault (Chapter~13).

\file{BADTOTAL.PRG} totals the ledger, but its programmer misspelled the field
\texttt{AMOUNT}:
\ProgramListing{BADTOTAL.PRG}

\ATEHead{ON ERROR}
\cmd{ON ERROR} names a command to be carried out whenever an error occurs.  While
the command runs, \fn{ERROR()} gives the number of the error.  A program uses it
to record the error before the message is displayed:
\begin{Steps}
\item Type \type{USE INVOICE}.
\item Type \type{ON ERROR STORE ERROR() TO MERR}.
\item Type \type{DO BADTOTAL}.
\item Type \type{? MERR, RECNO()}, then \type{ON ERROR} to remove the trap.
\end{Steps}
\DotExample{examples/p12-onerror}
The program stopped on its first record, with error 12.

\ATEHead{Where Do You Go From Here?}
Chapter~13 shows how to find and correct the mistake.

\ATEHead{Summary}
\begin{enumerate}
\item An error message gives a number and a text; Using InitechBase, Chapter~7,
  explains each.
\item A program stops at the command that fails; its memory variables remain.
\item \cmd{ON ERROR} \xp{command} runs a command when an error occurs;
  \fn{ERROR()} gives its number.
\end{enumerate}
